\documentclass[11pt]{article}

\usepackage[margin=1in]{geometry}
\usepackage{hyperref}
\usepackage{booktabs}
\usepackage{xcolor}
\usepackage{enumitem}
\usepackage{listings}
\usepackage{parskip}
\usepackage{caption}

\hypersetup{
    colorlinks=true,
    linkcolor=blue,
    urlcolor=blue,
    citecolor=blue
}

\lstset{
    basicstyle=\ttfamily\small,
    breaklines=true,
    frame=single,
    backgroundcolor=\color{gray!10},
    columns=fullflexible
}

\title{MCP-SecBench: A Defensive Proxy for Mitigating\\ Attacks Against MCP-Based AI Agents}
\author{Kavya Sharma}
\date{September 18, 2026}

\begin{document}

\maketitle

\begin{abstract}
The Model Context Protocol (MCP) lets AI agents connect to external tools --
file systems, search engines, terminals, databases -- turning a
conversational model into one that can act. This capability introduces a new
attack surface: a malicious or compromised MCP tool can attempt to hijack the
agent through its name, its parameter schema, or its response text. This
report describes an extended, original implementation built on top of the
\emph{MCP Security Bench} (MSB) benchmark~\cite{msb2026}, which measures how
often such attacks succeed (Attack Success Rate, ASR) across four attack
surfaces. Our contribution is \texttt{mcp\_guard}, a defensive proxy layer
that validates tool signatures, gates tool parameters, and sanitizes tool
responses without modifying MSB or its underlying agent framework, plus a
unified CLI tool for running baseline-vs-defended comparisons entirely on a
free, local, offline language model. We report results from a 120-scenario
batch run, discuss what the data does and does not show, and outline the
concrete next steps (a larger local model, full attack-surface coverage)
needed to turn this from a working first implementation into a complete
evaluation.
\end{abstract}

\tableofcontents

\section{Introduction: MCP and the Security Problem}

The Model Context Protocol (MCP) is a standard that lets AI agents invoke
external tools: reading files, running searches, controlling a terminal,
querying a database, and more. Where a purely conversational model can only
respond in text, an MCP-connected agent can \emph{act} on the user's behalf.

That capability is also the risk. An MCP tool is, functionally, a small
program the agent is told to trust. If that tool -- or content it returns --
is malicious, it can attempt to manipulate the agent into performing actions
the user never requested: exfiltrating private data, corrupting files,
terminating unrelated processes, or planting unauthorized access credentials.
This is a qualitatively different problem from classical software security:
the question is not whether an attacker can breach the agent's execution
environment, but whether the agent can be deceived by the very tools it was
instructed to trust.

MSB identifies four distinct surfaces through which this deception can occur,
summarized in Table~\ref{tab:surfaces}.

\begin{table}[h]
\centering
\caption{MSB's four MCP attack surfaces.}
\label{tab:surfaces}
\begin{tabular}{@{}p{3.3cm}p{10cm}@{}}
\toprule
\textbf{Attack surface} & \textbf{Mechanism} \\
\midrule
Tool Signature & A tool's name collides with a trusted tool, or its
description contains a hidden instruction (prompt injection, preference
manipulation). \\
Tool Parameter & The tool's schema requests an out-of-scope argument the
agent should not need to supply (e.g.\ an internal model identifier). \\
Tool Response & The result returned by a tool call is a fabricated error or
a fake ``urgent new instruction'' designed to hijack the agent mid-task. \\
Retrieval Injection & A retrieved document or search result embeds hidden
instructions disguised as ordinary content. \\
\bottomrule
\end{tabular}
\end{table}

\section{The MSB Benchmark}

MCP Security Bench (MSB)~\cite{msb2026} is an ICLR 2026 benchmark that
automates exactly this kind of test: it instantiates a real LLM agent
connected to real MCP tool servers, assigns it an everyday task (e.g.\
``search arXiv for papers on machine learning'', ``get the current working
directory''), and covertly injects one of twelve attack variants into the
pipeline. After execution, MSB checks whether the attack's concrete goal was
achieved -- for example, whether a fabricated SSH key was written to a file,
whether a private phone number was exfiltrated, or whether an unrelated
process was killed.

MSB's own published results, evaluated across ten LLM backbones, report an
average ASR of \textbf{40.35\%} across all attack types, with the
Out-of-Scope-Parameter attack alone averaging \textbf{76.50\%}. Critically,
MSB is purely a \emph{measurement} instrument: it contains no mitigation of
any kind. This is the gap the present work addresses.

\section{System Design}

We preserved \texttt{baseline/} as an unmodified copy of upstream MSB and
built four new components on top of it, described below and summarized
in Table~\ref{tab:contrib}.

\subsection{Defensive Proxy: \texttt{mcp\_guard.py}}

The core contribution is \texttt{mcp\_guard.py}, a middleware layer that
composes with an already-initialized MCP agent \emph{without} modifying the
agent framework's source code. It performs three checks, corresponding to
the three attack surfaces that are actually reachable at the agent-execution
layer (Retrieval Injection is discussed as a current limitation in
Section~\ref{sec:limitations}):

\begin{enumerate}[label=(\arabic*)]
\item \textbf{Tool Signature Validation.} Before the agent's system prompt is
constructed, the tool list is de-duplicated by name (defeating name-collision
attacks), and each tool's description is scanned for known injection-template
markers, which are stripped before the description reaches the prompt.
\item \textbf{Tool Parameter Gating.} Before a tool call is dispatched, its
arguments are checked against a set of disallowed keys (MSB's own
out-of-scope markers such as \texttt{llm\_model\_name}, plus generic
secret-like keys such as \texttt{api\_key} or \texttt{password}). High-risk
tool calls (\texttt{write\_file}, \texttt{edit\_file}, \texttt{kill\_process},
\ldots) receive additional content inspection for SSH/PEM key material or
references to sensitive filenames.
\item \textbf{Tool Response Sanitization.} After a tool call returns (or
after MSB's own attack-simulation logic substitutes a canned malicious
string), the observation text is scanned for MSB's known attack template
fragments and generic instruction-override phrasing, and redacted before it
is appended to the agent's context.
\end{enumerate}

Technically, this is implemented as a per-instance monkey-patch of the
LangChain \texttt{AgentExecutor}'s tool-dispatch method, applied only when
\texttt{apply\_defense()} is explicitly called. This design choice is what
makes the baseline-vs-defended comparison valid: the same scenario, same LLM,
and same random conversational trajectory are used in both conditions: the
only variable is whether the defensive wrapper is attached.

\subsection{Windows-Compatible Setup}

MSB's own \texttt{setup.py} assumes a POSIX shell (it shells out to
\texttt{which uv}) and destructively \texttt{shutil.move}s a LangChain patch
file into the installed package, which would silently delete it from the
tracked repository. \texttt{setup\_windows.py} resolves \texttt{uv} via
\texttt{shutil.which}, writes correctly JSON-escaped Windows paths into every
MCP server configuration file, and \emph{copies} rather than moves the patch,
keeping \texttt{baseline/} a complete and auditable mirror of upstream.

\subsection{Pluggable LLM Backend}

MSB hardcodes its model selection to paid providers (OpenRouter, DeepSeek,
Alibaba Tongyi). We generalized this to an explicit
\texttt{<backend>/<model>} specification supporting \texttt{ollama/},
\texttt{openrouter/}, \texttt{openai/}, and \texttt{deepseek/} backends. This
allowed the entire evaluation in this report to be run against a free,
fully local, offline model (Ollama), with zero API cost.

\subsection{Unified CLI Runner}

\texttt{cli\_runner.py} exposes three modes:

\begin{lstlisting}
python cli_runner.py --mode baseline --llm ollama/llama3.2:3b
python cli_runner.py --mode defended --llm ollama/llama3.2:3b
python cli_runner.py --mode compare
\end{lstlisting}

It reuses MSB's own scenario-construction helpers
(\texttt{tool\_exist}, \texttt{complete\_server\_config}) directly from
\texttt{baseline/main.py}, and implements a success oracle that reads the
sandboxed \texttt{operation\_space/output} artifact directly, rather than
regex-parsing MSB's own log text as \texttt{metrics.py} does.

\begin{table}[h]
\centering
\caption{Summary of original contributions vs.\ unmodified upstream content.}
\label{tab:contrib}
\begin{tabular}{@{}ll@{}}
\toprule
\textbf{Component} & \textbf{Status} \\
\midrule
\texttt{baseline/} & Unmodified copy of upstream MSB \\
\texttt{defenses/mcp\_guard.py} & Original -- defensive proxy middleware \\
\texttt{setup\_windows.py} & Original -- Windows-compatible setup \\
\texttt{cli\_runner.py} & Original -- unified baseline/defended/compare CLI \\
\texttt{run\_batch.py} & Original -- curated batch driver \\
\texttt{custom\_scenarios/} & Original -- extension scaffold (not yet populated) \\
\bottomrule
\end{tabular}
\end{table}

\section{Experimental Setup}

All experiments were run at zero monetary cost:

\begin{itemize}
\item \textbf{Model:} \texttt{llama3.2:3b}, a small open-weight model served
locally via Ollama.
\item \textbf{Tools:} MSB's real MCP servers -- a filesystem server, a
terminal-control server, and a web-search server -- each spawned as an
actual subprocess per scenario, not simulated.
\item \textbf{Scenario grid:} 3 attack types spanning three attack surfaces
(\texttt{prompt\_injection}, \texttt{out\_of\_scope\_parameter},
\texttt{false\_error}) $\times$ 5 canonical attack goals $\times$ 2 tools,
executed once undefended and once with \texttt{mcp\_guard} attached, for
120 total scenario runs.
\end{itemize}

Reaching a working local setup required resolving several practical
obstacles: adapting MSB's POSIX-only setup script for Windows, and three
failed attempts at installing Ollama (a hung installer requiring elevated
permissions with no interactive session to grant it, a corrupted archive
download, and a severely rate-limited connection) before a working local
model was obtained via a direct portable-archive download.

\section{Results}

Table~\ref{tab:results} reports Attack Success Rate (ASR) for the
undefended baseline and the \texttt{mcp\_guard}-defended condition across
the three tested attack types.

\begin{table}[h]
\centering
\caption{Attack Success Rate, baseline vs.\ defended (\texttt{llama3.2:3b}, 120 scenarios).}
\label{tab:results}
\begin{tabular}{@{}lccc@{}}
\toprule
\textbf{Attack type} & \textbf{Baseline ASR} & \textbf{Defended ASR} & \textbf{Reduction} \\
\midrule
Prompt Injection & 0.0\% & 0.0\% & 0.0 pp \\
Out-of-Scope Parameter & 0.0\% & 0.0\% & 0.0 pp \\
False Error & 0.0\% & 0.0\% & 0.0 pp \\
\bottomrule
\end{tabular}
\end{table}

The headline ASR is 0\% in both conditions, which requires careful
interpretation rather than a simple pass/fail reading.

\textbf{The underlying model did not fall for any tested attack, even
undefended.} A 3-billion-parameter local model may lack the multi-step
tool-orchestration reliability needed to carry an injected instruction
through to a concrete malicious action (e.g.\ actually calling
\texttt{write\_file} with attacker-supplied content), independent of whether
a defense is present. This leaves no observable gap for the ASR metric
itself to show a reduction against.

\textbf{The guard was nonetheless measurably active.} Across the 60
defended-mode scenarios, \texttt{mcp\_guard} recorded \textbf{31 distinct
interventions} (redacting injected tool-response content or refusing
out-of-scope parameters), averaging 0.52 interventions per scenario. The
corresponding count in the 60 baseline scenarios was exactly zero, by
construction (the guard is never attached). This is direct, quantitative
evidence that the mitigation layer is functioning as designed, independent
of the model's own resistance to the underlying attack.

\textbf{An unrelated infrastructure limitation affected 8 of 60 scenarios
per mode.} All eight are \texttt{kill\_process} attack-task scenarios, which
depend on MSB's Desktop Commander support tool; that tool is currently
fetched through a third-party hosted registry (Smithery) that returns a
``server not found'' error in this environment. The failure is identical and
symmetric across baseline and defended runs, so it does not bias the
comparison, but it does mean the \texttt{kill\_process} results are not
currently informative.

\section{Limitations and Future Work}
\label{sec:limitations}

\begin{itemize}
\item \textbf{Model capacity confound.} The single most important next step
is repeating this evaluation against a more capable model. A larger local
model (\texttt{llama3.1:8b}, matching a model MSB itself reports at
approximately 20\% average ASR) has already been downloaded and is queued
for this comparison; it should surface genuine baseline attack successes,
enabling a real before/after reduction figure rather than a 0\%-to-0\%
comparison.
\item \textbf{Retrieval Injection is not yet scored.} MSB's own success
criterion for this attack type is qualitatively different from the other
three (whether the agent was deceived into consulting the wrong document,
rather than whether a file artifact was planted), and the current success
oracle does not implement that logic. It is deliberately excluded from batch
runs rather than silently reported as a misleading 0\%.
\item \textbf{Third-party tool dependency.} \texttt{kill\_process} scenarios
require an alternative source for the Desktop Commander tool, independent of
Smithery's hosted registry.
\item \textbf{Scope.} This is a first working implementation covering three
of MSB's twelve attack-type combinations. The natural extension is full
coverage of MSB's attack grid and validation across multiple model sizes to
establish whether the observed guard-intervention rate scales with a model's
underlying susceptibility.
\end{itemize}

\section{Conclusion}

This work demonstrates that a defensive proxy layer can be attached to an
MCP-based agent pipeline without modifying the underlying benchmark or agent
framework, and that its interventions can be measured directly and
independently of whether the underlying attack succeeds. Our first
end-to-end evaluation, run entirely at zero cost on a local model, did not
find a model vulnerable enough to demonstrate an ASR reduction, but it did
establish -- with concrete counts, not just code review -- that the guard is
actively intercepting attack-shaped content. Extending this evaluation to a
more capable model and to MSB's full attack-type grid is the immediate next
step toward a complete assessment of \texttt{mcp\_guard}'s effectiveness.

\begin{thebibliography}{9}
\bibitem{msb2026}
Zhang, D., Li, Z., Luo, X., Liu, X., Li, P.\ P., Xu, W.
\emph{MCP Security Bench (MSB): Benchmarking Attacks Against Model Context
Protocol in LLM Agents.} The Fourteenth International Conference on Learning
Representations, 2026. \url{https://openreview.net/forum?id=irxxkFMrry}
\end{thebibliography}

\end{document}
